% Teaching example from topic 1.3: y^2 = x^3 + x + 1 modulo 17.
\begingroup
\setstretch{1}
\hyphenpenalty=10000
\newcommand{\ecblock}[3]{%
 {\fontsize{13}{15}\selectfont\bfseries\boldmath #1}\\[4pt]
 {\fontsize{10}{12}\selectfont #2}\\[3pt]
 {\fontsize{8}{10}\selectfont (#3)}%
}
\begin{tikzpicture}[x=1cm,y=1cm,
 box/.style={proceso,text width=30mm,minimum height=18mm,inner sep=6pt},
 size/.style={etiqueta,fill=none,font=\sffamily\fontsize{9}{11}\selectfont}]
 \node[font=\sffamily\fontsize{10}{12}\selectfont,align=center] at (1.85,5.2)
   {Small-field example\\$y^2\equiv x^3+x+1\pmod{17}$};
 \begin{scope}[x=0.22cm,y=0.22cm]
  \draw[step=1,draw=gray!20,line width=0.2pt] (0,0) grid (16,16);
  \draw[-{Latex[length=1.5mm]},draw=Grafito] (0,0)--(17.5,0) node[right,font=\scriptsize] {$x$};
  \draw[-{Latex[length=1.5mm]},draw=Grafito] (0,0)--(0,17.5) node[above,font=\scriptsize] {$y$};
  \foreach \t in {0,8,16}{
   \node[below,font=\sffamily\scriptsize] at (\t,0) {\t};
  }
  \foreach \t in {8,16}{\node[left,font=\sffamily\scriptsize] at (0,\t) {\t};}
  \foreach \x/\y in {0/1,0/16,4/1,4/16,6/6,6/11,9/5,9/12,10/5,10/12,11/0,13/1,13/16,15/5,15/12,16/4,16/13}{\fill[Grafito!60] (\x,\y) circle (1.3pt);}
  \fill[BordePublico] (6,6) circle (2.6pt);
  \node[above left,font=\sffamily\small,text=BordePublico] at (6,6) {$G$};
  \fill[BordePublico] (9,5) circle (2.6pt);
  \node[below right,font=\sffamily\small,text=BordePublico] at (9,5) {$Q_A$};
 \end{scope}
 \node[box,fill=white,funcionsecreta] (private) at (7.7,4.9)
   {\ecblock{$d_A=2$}{Private key}{secret scalar}};
 \node[box,funcioncripto] (multiply) at (7.7,1.8)
   {\ecblock{$Q_A=d_AG$}{Point multiplication}{$G+G$ in this example}};
 \node[box,fill=white,funcionpublica] (public) at (12.4,1.8)
   {\ecblock{$Q_A=(9,5)$}{Public key}{point on the same curve}};
 \coordinate (generator) at (1.32,1.32);
 \draw[flujo] (private.south) -- (multiply.north);
 \draw[flujo] (generator) -- (4.5,1.32) -- (4.5,1.8)
   -- node[size,above] {$G=(6,6)$} (multiply.west);
 \draw[flujo] (multiply.east) -- (public.west);
 \leyendaFunciones{(6.2,-1.3)}
\end{tikzpicture}
\endgroup
